\documentclass[11pt]{article}

\usepackage[margin=1in]{geometry}
\usepackage[T1]{fontenc}
\usepackage[utf8]{inputenc}
\usepackage{lmodern}
\usepackage{microtype}
\usepackage{amsmath,amssymb,amsthm}
\usepackage{booktabs}
\usepackage{array}
\usepackage{enumitem}
\usepackage{listings}
\usepackage[hidelinks]{hyperref}
\usepackage{xurl}

\newtheorem{definition}{Definition}
\newtheorem{theorem}{Theorem}
\newtheorem{remark}{Remark}

\newcommand{\eps}{\varepsilon}
\newcommand{\TV}{\mathrm{TV}}
\newcommand{\KL}{\mathrm{KL}}

\lstset{basicstyle=\footnotesize\ttfamily,breaklines=true,frame=single,columns=fullflexible}

\newcommand{\figbox}[1]{%
  \begin{center}
  \fbox{\parbox{0.93\linewidth}{\vspace{0.5em}\centering\emph{Figure omitted in the source-only archive.}\\[0.25em]\textbf{#1}\vspace{0.5em}}}
  \end{center}
}

\title{W-Congestion-EQ: Verifiable Congestion Budgets\\
Proof-Carrying Certificates and Transparency Receipts for Anonymous Overlays}
\author{}
\date{Draft v0.41 (2026-02-24)}

\begin{document}
\maketitle

\begin{abstract}
Congestion-EQ budgets are meaningful only if they are verifiable by parties other than the service operator.
We present a practical assurance layer for anonymous overlays (and Anonymous-DHT lookup services) that publishes compact \emph{congestion budget certificates} binding: (i) the declared congestion witness family, (ii) the PSC-Q (or other) compiler knobs that implement it, and (iii) the accountant arithmetic that yields an end-to-end budget.
To reduce trust in a single publisher, certificates are anchored in append-only transparency logs and distributed as \emph{receipts} with inclusion/consistency proofs.

The contribution of this paper is the verification surface and its threat model: what is mechanically checkable offline, what policies reduce equivocation (split-world), and how to connect congestion certificates to release-time change control.
We keep this paper brief by importing generic receipt/log semantics from the Release-Property paper and focusing on congestion-specific certificate fields, verifiers, and deployment policies.
\end{abstract}

\paragraph{Scope.}
We define the \emph{assurance layer} for congestion budgets: certificate fields, verifier checks, and anti-equivocation publication policies for Congestion-EQ style claims.

\paragraph{Units.} Budgets are published in bits (use $\log_2$ and $2^{\epsilon}$); results stated in nats can be converted by dividing by $\ln 2$ (see Synthesis~8).\cite{series:notation}

\paragraph{Imports.}
We import the congestion witness family from Congestion-EQ and mechanisms from PSC-Q.\cite{series:congeq,series:pscq}
Generic receipt/log semantics and release-time gating are treated as owned by the Release-Property workflow.\cite{series:release}
Packaging and versioning are treated as owned by ABOM/OINL and disagreement-gated recertification.\cite{series:abom,series:recert}

\paragraph{Exports.}
A \emph{W-Congestion-EQ receipt} schema: signed epoch-level congestion budgets binding (i) the witness family, (ii) declared knob sets, and (iii) accountant arithmetic; optionally anchored into transparency logs for split-world detection.

\paragraph{Artifact policy.}
Source-only; schemas, verifiers, and microstudy plots referenced below are available on request.\cite{series:artifactpolicy}

\section{Position in the series}
Congestion-EQ defines the \emph{witness} and \emph{budget metric} (TV constraints over congestion observables) \cite{series:congeq}.
PSC-Q defines a compiler-friendly mechanism family that meets these budgets under declared calendars and substitution padding \cite{series:pscq}.
This paper (W-Congestion-EQ) defines the \emph{assurance layer}: how a deployment turns ``we run PSC-Q with knob set $\theta$'' into a checkable, replayable statement that third parties can verify.

We treat receipt formats, transparency logs, and release-time gating as imported infrastructure from the release-property workflow \cite{series:release}.
ABOM/OINL is the packaging companion: a certificate should point to an ABOM entry identifying its schema versions and evidence digests \cite{series:abom}.
Change control (what happens when you update knobs) is governed by disagreement-gated recertification \cite{series:recert}.

\section{Threat model and goals}
\textbf{Participants.}
A provider runs an anonymity overlay and publishes periodic (e.g., daily) congestion-budget claims.
Users/auditors consume those claims.
\textbf{Adversary.}
We consider a provider who may (i) misreport knobs or accountant outputs, (ii) publish evidence selectively, or (iii) equivocate by presenting inconsistent certificates to different audiences.
\textbf{Goal.}
Define a certificate/receipt interface with an offline verifier that:
(1) checks internal arithmetic consistency; (2) binds claims to declared knobs and witness families; and (3) reduces equivocation via transparency receipts.

\textbf{Non-goals.}
A passing certificate is not a proof of real-world anonymity: it does not prevent hidden side channels.
It is a verifiable statement about what was claimed, how it was computed, and whether the claim is consistent with the public contract objects of the series.

\section{Contract objects and what must be bound}
\subsection{Witness binding (Congestion-EQ import)}
Congestion-EQ represents congestion leakage through a witness $W$ (queueing delay, RTT, loss/ECN marks, or derived aggregates), and budgets the TV distance between witness distributions under different destinations/queries \cite{series:congeq}.
A congestion certificate must bind:
\begin{itemize}[leftmargin=1.2em]
\item the witness family identifier (including units and sampling rule);
\item the policy for truncation/windowing (to avoid unbounded-support pathologies);
\item the budget statement (per epoch and composed) including any reset/flush model.
\end{itemize}

\subsection{Knob binding (PSC-Q import)}
PSC-Q implements congestion equalization with public calendars, substitution padding, and (optionally) declared dithers \cite{series:pscq}.
A congestion certificate must bind the deployed knob set $\theta$, including:
\begin{itemize}[leftmargin=1.2em]
\item calendar structure (period, slotting, any tiered calendars);
\item substitution padding policy (cover distribution; feasibility constraints);
\item retry discipline relevant to congestion observables (schedule-only retries, if any).
\end{itemize}

\subsection{Packaging hooks (ABOM/OINL import)}
We treat the certificate as an ABOM-referenced artifact:
\begin{itemize}[leftmargin=1.2em]
\item \textbf{Schema pinning.} The certificate names a schema version and canonicalization profile.
\item \textbf{Digest binding.} The certificate carries content hashes for any evidence objects referenced (e.g., topology snapshot digest, calibration profile digest) even if those objects are not published in the source-only archive.
\end{itemize}
The point is not to publish everything, but to let a verifier say ``I checked the claim under this exact set of inputs''.

\section{Certificate format (congestion-specific core)}
We model certificates as canonicalized JSON objects, signed, and optionally anchored in a transparency log.
Generic JSON canonicalization and receipt semantics are imported from \cite{series:release}; here we only specify the congestion-specific \emph{fields and checks}.

\subsection{Minimal field set}
A minimal congestion certificate contains:
\begin{itemize}[leftmargin=1.2em]
\item \texttt{id}: deployment identifier and epoch range;
\item \texttt{witness}: witness family id, sampling and truncation parameters;
\item \texttt{mechanism}: PSC-Q knob manifest $\theta$ (calendar + padding);
\item \texttt{accountant}: per-epoch budget terms and composed $\eps$ (or $(\eps,\delta)$ if used);
\item \texttt{evidence\_digests}: hashes for calibration profiles and reference data;
\item \texttt{signatures}: signer id and signature over canonical bytes;
\item \texttt{receipts} (optional): transparency log receipts for the signed bytes.
\end{itemize}

\subsection{Canonicalization}
JSON has many syntactic representations of the same object.
Therefore signing must be done over a canonical byte representation (e.g., RFC~8785/JCS) \cite{rfc8785}.
This is a MUST check: a verifier recomputes canonical bytes and verifies the signature over those bytes.

\section{Mechanical verification checklist}
We separate what a verifier can \emph{mechanically} check from what is an ecosystem policy choice.

\subsection{MUST checks}
\begin{enumerate}[leftmargin=1.6em]
\item \textbf{Schema and canonicalization.} Validate the JSON object against a pinned schema version; canonicalize using the declared profile (e.g., JCS) \cite{rfc8785}.
\item \textbf{Signature.} Verify the signature over canonical bytes; verify key identity according to the trust model used by the deployment (e.g., pinned keys, or an identity framework).
\item \textbf{Arithmetic.} Recompute the accountant from the declared per-epoch terms and repetition/reset rules; check that the composed bound matches the certificate.
\item \textbf{Binding.} Check that the witness id, knob manifest, and evidence digests are present and well-formed.
\end{enumerate}

\subsection{SHOULD checks (anti-equivocation)}
If transparency receipts are present:
\begin{enumerate}[leftmargin=1.6em]
\item \textbf{Inclusion.} Verify inclusion proof for the certificate hash against the signed log checkpoint.
\item \textbf{Consistency.} Verify consistency proofs between checkpoints to ensure an append-only view (monotone tree size) \cite{rfc9162}.
\item \textbf{Witness policy.} Require a witness/monitor policy (or cross-logging) that makes split-world attacks detectable at acceptable probability.
\end{enumerate}
The mechanics follow Certificate Transparency v2-style receipts \cite{rfc9162} or generic receipt formats (e.g., COSE receipts) \cite{cose_receipts_draft}.

\section{Split-world risk and a simple detectability law}
Transparency logs reduce equivocation only if clients/auditors obtain \emph{comparable views}.
A split-world operator could show different certificates to different audiences while maintaining two log views.
CT-style defenses rely on gossip, cross-logging, and monitors \cite{rfc9162}.

We use a minimal probabilistic model to guide policy choices.
Assume the operator equivocates on a fraction $\rho$ of epochs (publishing inconsistent certificates).
Let $m$ independent witnesses each sample $s$ epochs uniformly at random and gossip checkpoints pairwise.
Under a standard ``collision'' approximation, the probability that at least one witness observes an equivocated epoch is
\begin{equation}
\Pr[\text{detect}] \;\approx\; 1 - (1-\rho)^{ms}.
\end{equation}
This ignores timing/correlation and is optimistic, but it gives a sizing rule: to make $\Pr[\text{detect}]\ge 1-\eta$, it suffices that $ms \gtrsim \frac{\log_2(1/\eta)}{\rho}$.
In practice, witnesses should bias sampling toward high-risk periods (software updates, key rotations) and enforce monotone-checkpoint rules.

\section{Connecting to release-time change control}
Congestion certificates are not static: knobs and deployments change.
We recommend treating each certificate as a \emph{release artifact} governed by the release-property workflow \cite{series:release}:
\begin{itemize}[leftmargin=1.2em]
\item publish certificates through a policy gate that checks MUST conditions;
\item anchor releases in a claims ledger and map them to ABOM/OINL entries \cite{series:abom};
\item require recertification policies keyed to disagreement thresholds \cite{series:recert}.
\end{itemize}
This keeps ``privacy budgets'' from becoming marketing text: they become attestations with change control.

\section{Discussion and limitations}
\paragraph{What this does not prevent.}
A provider can satisfy certificate arithmetic while leaking through channels not modeled in the witness family.
The right interpretation is: ``the published budget is consistent with the public contract objects.''
\paragraph{Why this still helps.}
It turns a privacy claim into an object that can be audited, diffed, and monitored over time (like SBOM-style supply-chain artifacts, but for anonymity budgets) \cite{torresarias_intoto_sec19,sigstore_acm2022,slsa_v1}.

\section{Takeaways}
\begin{itemize}[leftmargin=*]
\item \textbf{W-Congestion-EQ is verifiability:} it turns congestion budgets into certificates that third parties can check without trusting the compiler.
\item \textbf{Receipts beat narratives:} publish a verifier report plus inclusion/consistency proofs in an append-only log to reduce equivocation and enable monitoring.
\item \textbf{Interfaces are small:} the certificate binds (i) a concrete plan/manifest and (ii) a dual witness that upper-bounds the congestion objective under the Congestion-EQ accountant.
\end{itemize}

\begin{thebibliography}{99}\small
\bibitem{series:congeq}
Anonymous.
\newblock \emph{Congestion-EQ: Total-Variation Budgets for Congestion Witnesses in Anonymous Overlays}.
\newblock Congestion series Paper~1, The-One-True-Archive (2026).

\bibitem{series:pscq}
Anonymous.
\newblock \emph{PSC-Q: Public Service Calendars and Substitution Padding for Congestion-Equalized Anonymous Lookups}.
\newblock Congestion series Paper~2, The-One-True-Archive (2026).

\bibitem{series:release}
Anonymous.
\newblock \emph{Privacy as a Release Property: Proof-Carrying Anonymity Receipts for Anonymous DHTs}.
\newblock Release \& destination Paper~A, The-One-True-Archive (2026).

\bibitem{series:abom}
Anonymous.
\newblock \emph{ABOM and OINL: An Anonymity Bill of Materials and an Odds-Inflation Nutrition Label}.
\newblock Anonymity series Paper~C, The-One-True-Archive (2026).

\bibitem{series:recert}
Anonymous.
\newblock \emph{Disagreement-Gated Recertification: Change Control for Certified Anonymous Systems}.
\newblock Certified series Paper~C, The-One-True-Archive (2026).

\bibitem{series:artifactpolicy}
Anonymous.
\newblock \emph{Artifact and Disclosure Policy for the One-True-Archive (Source-Only Public Release)}.
\newblock Synthesis~6, The-One-True-Archive (2026).

\bibitem{rfc8785}
A.~Rundgren.
\newblock {RFC}~8785: {JSON} Canonicalization Scheme ({JCS}).
\newblock IETF, 2020.

\bibitem{rfc9162}
B.~Laurie, E.~Kasper, E.~Messeri, and R.~Stradling.
\newblock {RFC}~9162: Certificate Transparency Version~2.0.
\newblock IETF, 2021.

\bibitem{torresarias_intoto_sec19}
S.~Torres-Arias, H.~Afzali, T.~K.~Kuppusamy, R.~Curtmola, and J.~Cappos.
\newblock in-toto: Providing farm-to-table guarantees for bits and bytes.
\newblock In \emph{USENIX Security}, 2019.

\bibitem{sigstore_acm2022}
Z.~Newman, J.~Speed Meyers, and S.~Torres-Arias.
\newblock Sigstore: Software signing for everybody.
\newblock In \emph{ACM CCS}, 2022. DOI: 10.1145/3548606.3560596.

\bibitem{slsa_v1}
SLSA contributors.
\newblock \emph{SLSA Specification v1.0}, 2023.

\bibitem{cose_receipts_draft}
IETF COSE Working Group.
\newblock \emph{COSE Receipts (Internet-Draft)}, 2025.

\bibitem{series:notation}
Anonymous.
\newblock Notation and Minimal Model for the One-True-Archive.
\newblock Series synthesis note (Synthesis~8), 2026. (This archive.)

\end{thebibliography}

\end{document}
